\scanpage{48}
\deleted{Si par contre $i$ est \unsure{égal|égale} à \unsure{$\infty$|$0$|$1$},}

En tous cas, si $i=0$, la situation est décrite par le fait que l'on a \deleted{fait} une réduction du groupe structural de $\mathfrak S_3$ à $\mathbb Z/2$ \inserted{[$X''$ \inserted{\unsure{donc|ou} $X'$} a une section !]}, et que tout est déterminé par la donnée d'un élément de $H^1(S,\mathbb Z/2)$. \deleted{Si par contre} \inserted{Et si} $i=1$, alors la situation est décrite en disant que l'on a une réduction du groupe structural : $\mathbb Z/3\subset\mathfrak S_3$ [\unsure{qui =|équiv.|grâce à} : une section de \inserted{la bisection} $X_1$] et la situation est décrite par la donnée d'un élément de $H^1(S,\mathbb Z/3)$.

\noindent\rule{.22\linewidth}{.4pt}

\deleted{\illegible{1}} En tous les cas encore, si $i$ est inversible, alors faisant opérer $\mathfrak S_3$ dans $P^1_{\deleted{\unsure{T}}S}/(\mathbb Z/3)$ par passage au quotient, on trouve dans \deleted{\unsure{ce dernier}} \inserted{\unsure{ce fibré croisé|le fibré associé|le fibré corresp.}} \deleted{\unsure{deux}} \uline{trois} \uline{sections} distinctes, \inserted{(prov\textsuperscript{nt} de $X',\unsure{X'''|X''}$ et $s$)} qu'on peut choisir comme $\infty,0,1$, d'où une identification de ce dernier à $P^1_{\deleted{\unsure{T|S}}S}$ ; et $X\longrightarrow P^1_S$ se factorise par
\diagram{\begin{tikzcd}[ampersand replacement=\&,column sep=1.6em,row sep=.25em]
X\arrow[r]\&P^1_S\arrow[r]\&P^1_S\\
\&\infty,0,1\arrow[r]\&\infty,0,i
\end{tikzcd}}
\[
t\mathbin{\unsure{-|\mapsto}}it^2
\]
; la bisection naturelle en question s'identifie alors \unsure{à|:} $\sqrt{i}$,
% archive sheet number: 47
